% Appendix, from exam/paper_b.md: Paper B, questions only.

\chapter{Self-Assessment Paper B}
\label{exb:ch:paper}

\noindent\textbf{Time:} 3 hours. \textbf{Closed book.} No compiler and no reference material.
\textbf{Total: 100 marks.}

\medskip
This paper exists to test your own skills and knowledge. It is not a qualification and it gates
nothing. It draws on all six chapters, so \textbf{it is meant to be taken once you have finished the
book}. The model answers are in Appendix~\ref{ansb:ch:answers}.

\section*{Rubric}
\textbf{C++17} is assumed throughout. Standard headers may be assumed included unless a question
asks you to name them.

\textbf{Code is marked on semantics, not on syntax.} A missing semicolon, a forgotten
\code{#include} or a brace in the wrong column costs nothing. A missing \code{virtual} on a
destructor costs everything.

\textbf{Where a question says ``in the book's style''}, it means the conventions used in the
chapters:
\begin{itemize}
  \item \code{camelCase} for functions and methods, type names with a leading capital.
  \item Private member variables prefixed \code{my}, private static members prefixed \code{our}.
  \item \code{static constexpr} members named with a leading capital, e.g.\ \code{MaxInstances}.
  \item Brace initialization, \code{noexcept} on driver methods, \code{[[nodiscard]]} on queries and
    repeated on overriding methods.
  \item Copy and move operations deleted on a class that represents a unique hardware resource.
\end{itemize}

\textbf{Where a question says ``state the consequence''}, naming the defect earns half the marks and
saying what it does to the running program earns the other half. ``This is wrong'' scores nothing.

\textbf{Where a question asks you to write a class}, documentation comments are not required and
are not marked. Write the code.

\begin{narrowtable}{@{}lcccccccc@{}}
  \thead{Question} & 1 & 2 & 3 & 4 & 5 & 6 & 7 & 8 \\
  \midrule
  \thead{Marks} & 12 & 12 & 12 & 13 & 13 & 13 & 12 & 13 \\
\end{narrowtable}

% ------------------------------------------------------------------------------------------
\question{The vocabulary}{12}
\label{exb:q:1}

\qpart{a} State what difference the extension \code{.h} or \code{.hpp} makes to the compiler, and
what actually decides whether a file is compiled as C or as C++. State what each extension signals
by convention, and say which of the two a vendor's register-definition header should carry, with the
reason. \worth{3}

\qpart{b} State what an anonymous namespace does and name the C construct it replaces. Then state
the purpose of a nested namespace such as \code{driver::gpio}, and write it in both the pre-C++17
form and the compact modern form. \worth{3}

\qpart{c} Give three reasons references are usually preferred over pointers when a function must
modify its argument. Then name one thing a pointer can do that a reference cannot, and give an
example from the book where that matters. \worth{3}

\qpart{d} Complete the function template below so that it sets every bit in the pack, using a C++17
fold expression rather than a loop. State what the compiler generates for the call
\code{set(reg, 1U, 2U, 3U)}. \worth{3}

\begin{cppcode}
template<typename T, typename... Bits>
constexpr void set(T& reg, const Bits... bits) noexcept
{
    static_assert(std::is_integral<T>::value,
        "Failed to set bit in register: T must be of integral type!");
    // Your code here.
}
\end{cppcode}

% ------------------------------------------------------------------------------------------
\question{The same driver, twice}{12}
\label{exb:q:2}
Below is a PWM driver as it would be written in C.

\begin{ccode}
#include <stdbool.h>
#include <stddef.h>
#include <stdint.h>

typedef struct
{
    uint8_t pin;
    uint16_t period_ms;
    bool enabled;
} pwm_t;

void pwm_init(pwm_t* self, const uint8_t pin, const uint16_t period_ms)
{
    if (NULL == self) { return; }
    self->pin       = pin;
    self->period_ms = period_ms;
    self->enabled   = false;
}

void pwm_set_enabled(pwm_t* self, const bool enabled)
{
    if (NULL == self) { return; }
    self->enabled = enabled;
}

bool pwm_is_enabled(const pwm_t* self)
{
    return NULL != self ? self->enabled : false;
}
\end{ccode}

\qpart{a} Rewrite it as a single-header C++ driver in the book's style. It shall:
\begin{itemize}
  \item live in namespace \code{driver::pwm}, as a class named \code{Pwm} that cannot be inherited
    from;
  \item replace \code{pwm_init} with a constructor that cannot be used for implicit conversions and
    whose period defaults to \code{20U};
  \item keep the pin and the period unchangeable after construction;
  \item expose \code{pin()}, \code{period_ms()}, \code{isEnabled()} and \code{setEnabled()}, with the
    query methods marked so that ignoring their result warns;
  \item be impossible to copy or move.
\end{itemize}
\noindent\worth{6}

\qpart{b} Name \textbf{four} distinct failure modes that exist in the C version and cannot occur in
your C++ version, and state for each what removes it. \worth{4}

\qpart{c} Your header and the vendor's C register header sit in the same project. State which
extension each should carry, and what must be added to the vendor's header before your driver can
call into it. \worth{2}

% ------------------------------------------------------------------------------------------
\question{Enumerations, files, and lifetime}{12}
\label{exb:q:3}

\qpart{a} Write the enumeration class \code{Direction} from Chapter~2 in namespace
\code{driver::gpio}, with the enumerators \code{Input}, \code{InputPullup} and \code{Output}, a fixed
underlying type, and whatever is needed to support the validity check below. Then write that check.

\begin{cppcode}
[[nodiscard]] constexpr bool isDirectionValid(Direction direction) noexcept;
\end{cppcode}

\noindent Give three reasons to prefer an enumeration class over a plain \code{enum}, and state what
fixing the underlying type buys on an embedded target. \worth{5}

\qpart{b} A class is split into \file{driver/gpio/gpio.hpp} and \file{driver/gpio/gpio.cpp}. For
each of the following, state whether it belongs in the header, in the source file, or in both, and
why:

\code{explicit}, a default argument, \code{[[nodiscard]]}, the trailing \code{const} on a query
method, \code{noexcept}, \code{= default} and \code{= delete}. \worth{4}

\qpart{c} State when the constructor runs and when the destructor runs for each of:
\begin{enumerate}
  \item a local object inside a function;
  \item an object at namespace scope;
  \item an object created with \code{new}.
\end{enumerate}
For the second, name the hazard this creates in an embedded program. \worth{3}

% ------------------------------------------------------------------------------------------
\question{What the compiler writes for you}{13}
\label{exb:q:4}

\qpart{a} The \code{Gpio} class from Chapter~2 declares its pin as \code{const std::uint8_t myPin;}.
A copy \emph{constructor} that makes a complete copy can be written for this class; a copy
\emph{assignment operator} that does so cannot. Explain why, and state what the compiler does about
the copy assignment operator when you do not declare one. \worth{4}

\qpart{b} A class \code{Buffer} owns a heap block through two members:

\begin{cppcode}
std::uint8_t* myData;
std::size_t mySize;
\end{cppcode}

\noindent Write its move constructor. Then state the one thing a move \emph{assignment operator} must
do that a move \emph{constructor} need not, and explain why the constructor is exempt. \worth{4}

\qpart{c} State what \code{= default} and \code{= delete} each mean. Give the block of operations a
driver class representing a unique hardware resource should delete, and the design reason.

Then answer this: writing \code{~Gpio() noexcept = default;} is often described as costing nothing,
because the compiler would have generated the same destructor anyway. State what it nevertheless
changes about the rest of the class. \worth{5}

% ------------------------------------------------------------------------------------------
\question{Designing against an abstraction}{13}
\label{exb:q:5}

\qpart{a} Write the interface \code{driver::serial::Interface} in the book's style. It shall support
initializing the port with a baud rate, writing a null-terminated string, and querying whether the
port has been initialized.

Then state, for four of the conventions \secref{c3:app:b} applies to an interface, what each is and
why it is there. \worth{5}

\qpart{b} State the difference between \code{private} and \code{protected} members. Name the three
kinds of inheritance in C++, state what each does to the access level of inherited members, and
state which one a concrete driver uses to implement an interface, and why. \worth{4}

\qpart{c} A colleague writes:

\begin{cppcode}
void sendGreeting(driver::serial::Interface serial) noexcept
{
    serial.write("Hello!\n");
}
\end{cppcode}

\noindent State what happens when this is compiled, and why. Then state what \emph{would} have
happened had the base class not been abstract, and name the mechanism. Give the corrected signature,
and state one reason not to make it \code{const}. \worth{4}

% ------------------------------------------------------------------------------------------
\question{Factories}{13}
\label{exb:q:6}

\qpart{a} Draw or describe the dependency graph of the four layers in the Chapter~4 example: system
logic, factory interface, concrete factories, concrete drivers. State which dependency is the one
that makes the system testable without hardware, and why. \worth{3}

\qpart{b} Write the stub factory \code{driver::factory::Stub} for the smart-pointer version of the
factory interface

\begin{cppcode}
[[nodiscard]] virtual std::unique_ptr<gpio::Interface> gpio(std::uint8_t pin) noexcept = 0;
\end{cppcode}

\noindent including whatever is needed to keep the compiler quiet about the unused pin. \worth{4}

\qpart{c} Give four differences between a raw \code{gpio::Interface*} and a
\code{std::unique_ptr<gpio::Interface>} as the return type of \code{gpio()}. State where
\code{std::move} becomes necessary once the logic class stores one, and what the source pointer holds
afterwards. \worth{4}

\qpart{d} Name one cost the factory pattern imposes on an embedded system, and describe the
alternative shown in Chapter~4 that avoids it entirely. \worth{2}

% ------------------------------------------------------------------------------------------
\question{A container with its size in the type}{12}
\label{exb:q:7}

\qpart{a} Write a class template \code{container::Array<T, Size>} in namespace \code{container}, in a
single header, providing:
\begin{itemize}
  \item a default constructor;
  \item \code{push(const T& element)}, returning \code{false} if the array is already full;
  \item \code{size()} and \code{capacity()};
  \item \code{operator[]} in both a modifiable and a read-only form;
  \item a compile-time check rejecting a size of zero, with a message.
\end{itemize}
The class shall not allocate memory. \worth{6}

\qpart{b} State why this class cannot have its method definitions placed in \file{array.cpp}, and
what the build would report if you tried. Then state how many distinct types the following three
declarations produce, with the reason, and what that means for the size of the binary. \worth{4}

\begin{cppcode}
container::Array<std::uint8_t, 8U> a{};
container::Array<std::uint8_t, 16U> b{};
container::Array<std::uint16_t, 8U> c{};
\end{cppcode}

\qpart{c} State one advantage \code{container::Array<T, Size>} has over the heap-based
\code{container::Vector<T>} from \secref{c5:app:c} on a target where dynamic allocation is
forbidden, and one thing you give up by putting the capacity in the type. \worth{2}

% ------------------------------------------------------------------------------------------
\question{Sharing state between threads}{13}
\label{exb:q:8}

\qpart{a} State the three conditions that together constitute a data race. State what the standard
says the program then does, and why that is a stronger statement than ``the value may be wrong''.
\worth{3}

\qpart{b} The TX and RX threads in Chapter~6 share this structure under a mutex:

\begin{cppcode}
struct SharedMem
{
    std::uint16_t data{};
    bool newData{false};
};
\end{cppcode}

\noindent A colleague deletes the mutex and writes instead:

\begin{cppcode}
struct SharedMem
{
    std::atomic<std::uint16_t> data{};
    std::atomic<bool> newData{false};
};
\end{cppcode}

\noindent State whether the data race is gone. State whether the program is correct, and justify
your answer with a specific interleaving. Then state the two things a mutex provides that a pair of
atomics does not. \worth{4}

\qpart{c} Name the two launch policies of \code{std::async} and state what each does. State what
happens when no policy is passed, and describe the failure this produces in code written as ``start
the work, do something else, collect the result later''. \worth{3}

\qpart{d} Describe priority inversion as a sequence of events involving three threads. State why
\code{std::mutex} cannot prevent it, and name what does. \worth{3}
